\documentclass[11pt]{article}
\usepackage{amsmath,amsthm,amssymb}
\usepackage[margin=1.05in]{geometry}
\usepackage{booktabs,array,longtable}
\usepackage[hidelinks]{hyperref}
\hypersetup{
  pdftitle={An Effective Classification of Congruence-Restricted Quadratic Relative-Error Records},
  pdfauthor={Elias De Jesus},
  pdfsubject={Diophantine approximation; quadratic irrationals; record sequences},
  pdfkeywords={Diophantine approximation, quadratic irrationals, record sequences,
    congruence conditions, binary quadratic forms, Pell equation, automorph group,
    reduction theory, relative error, exact arithmetic, certificates, effective algorithms}
}
\usepackage{microtype}

\newtheorem{theorem}{Theorem}[section]
\newtheorem{proposition}[theorem]{Proposition}
\newtheorem{lemma}[theorem]{Lemma}
\newtheorem{corollary}[theorem]{Corollary}
\theoremstyle{definition}
\newtheorem{definition}[theorem]{Definition}
\newtheorem{algorithm}[theorem]{Algorithm}
\theoremstyle{remark}
\newtheorem{remark}[theorem]{Remark}
\newtheorem{warning}[theorem]{Warning}
\newtheorem{example}[theorem]{Example}

\newcommand{\Z}{\mathbb{Z}}
\newcommand{\Q}{\mathbb{Q}}
\newcommand{\R}{\mathbb{R}}
\newcommand{\adm}{\mathcal{A}}
\newcommand{\dd}{\delta}
\newcommand{\epsbar}{\overline{\varepsilon}}
\newcommand{\epsw}{\varepsilon_w}
\newcommand{\sgn}{\operatorname{sgn}}
\DeclareMathOperator{\Nm}{N}

\title{An Effective Classification of Congruence-Restricted\\
Quadratic Relative-Error Records}
\author{Elias De Jes\'us\\
\small Independent Researcher\\
\small ORCID: \href{https://orcid.org/0009-0007-0190-9143}{0009-0007-0190-9143}\\
\small \texttt{dejesuselias10@gmail.com}}
\date{16 November 2026}

\begin{document}
\maketitle

\begin{abstract}
We give a terminating exact-arithmetic algorithm, with independently checkable certificates,
for classifying strict quadratic relative-error records under arbitrary congruence
restrictions; the classification includes imprimitive states, and the positive-conjugate cases
that require explicit backward exclusions.

Let $\alpha>0$ be a real quadratic irrational, given as either root of a primitive integral
quadratic $Ax^2+Bx+C$ with $A>0$ and nonsquare discriminant $\Delta=B^2-4AC$. For pairs of
positive integers $(m,k)$ whose residue $(m,k)\bmod M$ lies in a prescribed nonempty set
$\mathcal R\subseteq(\Z/M\Z)^2$, order the pairs by the weighted height $H=m\alpha+k$ and
measure them by the \emph{relative} error $\varepsilon=|m\alpha-k|/(m\alpha+k)$. A pair is a
\emph{strict record} if its error is smaller than at every admissible pair of smaller height.
We prove that the strict records always form a finite, explicitly decided part together with
finitely many certified forward orbits of one integral matrix, each starting at a computed
level, and we give a terminating algorithm that outputs this data with an exact certificate.
No coprimality restriction is imposed on the admissible pairs, and no sign or monicity
hypothesis on the polynomial: $\alpha$ may be the
smaller root, its conjugate may be positive, $\alpha$ may be less than $1$, $A$ may exceed $1$,
$\mathcal R$ may be any union of residue classes, and imprimitive records occur and are
classified. The mechanism is the proper automorph of the integral defect form
$Q=Ak^2+Bmk+Cm^2$: a power $P$ of it fixes every residue class modulo $M$ and preserves $Q$, an
explicit cutoff bounds $|Q|$ on records, classical reduction theory makes the surviving orbits
finite, and finitely many exact comparisons decide each orbit's tail. Two phenomena absent from
the monic negative-conjugate case require separate treatment and are isolated: when the
conjugate is positive the positive cone is not forward invariant and the height is not monotone
along an orbit, and the backward exclusion then splits into two cases with different
certificates. A separate negative result is isolated: a strict record need not be a vertex of
the convex hull of the admissible points of its cone, so the records cannot in general be read
off a Klein--Arnold sail; the witness is exact. That all sufficiently late records lie in
finitely many automorph orbits is classical once the defect is bounded; what is offered here as
a \emph{proposed} contribution is the effectiveness, the generality, and the completeness of the
certificates. Computations are used only to instantiate and audit the proof, not to establish
it; all arithmetic is exact. The note is part of an ongoing programme, and independent
verification, corrections and counterexamples are welcome.
\end{abstract}

\section{Introduction}\label{sec:intro}

Fix an irrational $\alpha>0$. Among pairs of positive integers $(m,k)$, how well can $k/m$
approximate $\alpha$ relative to the size of the pair, and which pairs are the successive
record holders? If ``size'' is the denominator $m$ and ``error'' is $|m\alpha-k|$, the answer is
classical: the records are the convergents of the continued fraction of $\alpha$. This note
concerns a different pairing, which arises when both coordinates carry weight:
\[
H(m,k)=m\alpha+k,\qquad
\varepsilon(m,k)=\frac{|m\alpha-k|}{m\alpha+k},
\]
so that the ordering quantity and the normalisation are the same. Two features make the
resulting record problem genuinely different from the classical one.

First, $\varepsilon$ is a monotone function of the \emph{multiplicative} error:
$\varepsilon=\bigl|\tanh\bigl(\tfrac12\log(k/(m\alpha))\bigr)\bigr|$, so records minimise
$|\log(k/m)-\log\alpha|$ under the weighted height. Proposition~\ref{prop:classical} compares
this with the \emph{absolute}-error criterion \emph{under the same height order} and gives the
one implication that holds, with a witness for its strictness. It says nothing about the
denominator-ordered notion of classical and parity-restricted best approximation: that is a
different ordering and no bridge to it is established here (Remark~\ref{rem:order2}). What is
established, and is of a geometric kind, is that records of this objective need not be
vertices of the convex hull of their own residue coset
(Proposition~\ref{prop:sail}) --- so the natural ``read the records off a sail'' reduction
fails.

Second, we restrict the admissible pairs by a congruence: $(m,k)\bmod M$ must lie in a
prescribed nonempty $\mathcal R\subseteq(\Z/M\Z)^2$. This is not a cosmetic restriction. An
arbitrary $\mathcal R$ is a \emph{union of cosets} of $M\Z^2$, not a sublattice, so one cannot
simply rescale; the gcd of an admissible pair need not be reducible inside the class, so
\emph{imprimitive records occur}; and the automorph group of the defect form must be replaced
by the finite-index subgroup fixing the classes.

The record set is infinite, and the question is whether it is \emph{finitely describable and
computable}. Theorem~\ref{thm:main} says it is, for every admissible input. The record set
always consists of an explicitly bounded finite region, decided directly, together with
finitely many complete forward orbits of one integral matrix $P$, each rebased at a computed
starting level. The proof is a chain of elementary steps around one invariant, the integral
defect
\[
Q(m,k)=Ak^2+Bmk+Cm^2,
\]
which $P$ preserves exactly; the finiteness comes from the classical reduction theory of
indefinite binary quadratic forms.

The work generalises three worked classifications in the author's earlier note
\cite{DeJesus2026oddodd} --- for $\alpha=\varphi$, $1+\sqrt3$ and $\sqrt3$ with $m,k$ both odd
--- which were proved one slope at a time under the hypotheses $B\le0$ (``$t\ge0$''),
$\alpha'<0$, $\alpha>1$, $A=1$, and the specific odd/odd class. All five hypotheses are removed
here. Two places where removing them is not routine are worth announcing, because they are
where the monic negative-conjugate intuition fails:

\begin{itemize}\itemsep2pt
\item When $\alpha'>0$ \emph{no power of the automorph has non-negative entries}, the positive
cone is not forward invariant, and an orbit of negative defect is positive in \emph{both}
directions (Lemma~\ref{lem:window}). Along such an orbit the height is U-shaped in the level,
so ``later'' does not mean ``higher''.
\item The natural error floor $\varepsilon\ge\sqrt\Delta/|B|$, which excludes the backward
portions in the negative-conjugate case, is \emph{false} on those doubly-infinite orbits: their
backward rays approach the floor from below. The exclusion therefore splits into two cases with
different certificates (Lemma~\ref{lem:backward}), one of them an explicit level cutoff
(Lemma~\ref{lem:cutlevel}).
\end{itemize}

\paragraph{What is classical, and what is proposed here.}
It is worth separating two statements that are easily conflated. The first is:

\begin{quote}
(a) all sufficiently late records lie in finitely many orbits of the automorph group of $Q$.
\end{quote}

This follows from the classical reduction theory of indefinite binary quadratic forms
\cite{BuellBQF,Cohen1993,ConradGenPell} once the defect is bounded --- Theorem~\ref{thm:cutoff}
supplies the bound, and Proposition~\ref{prop:box} is then a finiteness statement of a thoroughly
standard kind. We claim no novelty for it, and the Pell solution, the proper automorph, the
$Q$-invariance and the reduction box are all classical inputs used as such. The second
statement is:

\begin{quote}
(b) \emph{which} orbit tails are records; at \emph{which} level each one begins; and a finite
certificate that the finitely many exceptions have been completely enumerated.
\end{quote}

Statement (b) does not follow from (a). It is what the comparison formulas of
Section~\ref{sec:compare}, the finite ledger with both of its endpoints justified
(Section~\ref{sec:ledger}), the binding-state sufficiency lemma
(Lemma~\ref{lem:binding}) and the disjoint partition of Section~\ref{sec:partition} are for, and
it is what makes the output of Theorem~\ref{thm:main} an answer rather than a reduction.

The \emph{proposed} contribution of this note is therefore threefold: effectiveness --- a
terminating algorithm with exact arithmetic that outputs the record set itself, not a bound on
where it stabilises; generality --- arbitrary $(A,B,C)$, either root, $\alpha'>0$, $\alpha<1$,
arbitrary $\mathcal R$ and imprimitive states, all of which are excluded by the hypotheses of
\cite{DeJesus2026oddodd}; and certificate completeness --- an explicit, independently
checkable list of the quantities that a verifier needs in order to confirm that nothing was
missed. ``Proposed'' is meant literally: these are three claims awaiting a fuller comparison
with the literature, not established firsts. Congruence-restricted approximation has been studied from several
directions that the author has not exhausted: parity and odd/odd continued fractions
\cite{KimLeeLiao2022,KimLeeLiao2025}, rationals with restricted numerator and denominator
\cite{Elsner2000}, best approximation in a prime-power congruence setting
\cite{KushwahaSarma}, metric and counting results for approximation with congruence conditions
on the pair \cite{HanLee2025}, and the geometry of sails and rivers attached to indefinite
binary forms \cite{Karpenkov2013,SpaldingVeselov}. None of them treats the present objective
under the present order, as far as the author can tell; the specific distinctions are set out in
Section~\ref{sec:limits}, item~6. The absence of a result in the works consulted is not a
priority claim, and corrections of this reading are explicitly invited.

Computation plays a fixed and limited role. The theorem and every lemma are proved for
arbitrary admissible input; the computations instantiate the algorithm on $26$ inputs, emit a
machine-readable certificate for each, and audit those certificates with a separately written
checker. Section~\ref{sec:repro} states exactly what was run and what it does and does not
establish. All arithmetic is exact: elements of $\Q(\sqrt\Delta)$ are pairs of rationals and
every sign is decided by an integer comparison.

\paragraph{Open engagement.}
This technical note forms part of an ongoing research program. Independent verification,
corrections, related references, and suggestions for extensions are welcome. In particular,
feedback on the universal scope of the classification, the certificate completeness conditions,
and connections with congruence-restricted approximation theory would be valuable.

\section{Setting}\label{sec:setting}

Let $A,B,C$ be coprime integers with $A>0$ and $\Delta=B^2-4AC>0$ not a perfect square. Fix
$s\in\{+1,-1\}$ and set
\[
\alpha=\frac{-B+s\sqrt\Delta}{2A},\qquad
\alpha'=\frac{-B-s\sqrt\Delta}{2A},\qquad
\dd=\alpha-\alpha'=\frac{s\sqrt\Delta}{A},
\]
and require only $\alpha>0$. Thus $\alpha$ may be either root, $\alpha'$ may have either sign,
$\alpha$ may be smaller or larger than $1$, and $\alpha$ need not be an algebraic integer. Note
$A|\dd|=\sqrt\Delta$.

Fix $M\ge1$ and a nonempty $\mathcal R\subseteq(\Z/M\Z)^2$. The set of \emph{admissible states}
is
\[
\adm=\bigl\{(m,k)\in\Z_{>0}^2:\ (m\bmod M,\,k\bmod M)\in\mathcal R\bigr\}.
\]
\emph{No coprimality is imposed.} For a state put
\[
H=m\alpha+k>0,\quad
\xi=k-m\alpha,\quad
\eta=k-m\alpha',\quad
Q=Ak^2+Bmk+Cm^2=A\,\xi\eta,\quad
\varepsilon=\frac{|\xi|}{H}.
\]
$Q$ is the integral defect; $Q\ne0$ because $\alpha$ is irrational. Also set
\[
G=2\alpha A^2\eta^2+BQ,\qquad Y=A^2\eta^2 .
\]

\begin{lemma}[Basic identities]\label{lem:ident}
For every state,
\begin{equation}\label{eq:ident}
H=\frac{G}{A^2\eta\,\dd},\qquad
G=A^2\eta\,\dd\,H,\qquad
\varepsilon=\frac{|Q|\sqrt\Delta}{|G|},\qquad
H=\frac{2\alpha\eta}{\dd}\Bigl(1+u\Bigr),\quad u:=\frac{BQ}{2\alpha Y}.
\end{equation}
In particular $\sgn G=\sgn(\eta\,\dd)$, and with $z=\xi/H$ (so $|z|=\varepsilon$),
\begin{equation}\label{eq:Qz}
Q=\frac{H^2z\,(A\dd-Bz)}{2\alpha},\qquad\text{hence}\qquad
|Q|\le\frac{H^2\varepsilon\bigl(\sqrt\Delta+|B|\varepsilon\bigr)}{2\alpha}.
\end{equation}
\end{lemma}

\begin{proof}
From $\xi=k-m\alpha$ and $\eta=k-m\alpha'$ we get $\eta-\xi=m\dd$ and
$\alpha\eta-\alpha'\xi=k\dd$, so $m=(\eta-\xi)/\dd$, $k=(\alpha\eta-\alpha'\xi)/\dd$ and
\[
H=m\alpha+k=\frac{2\alpha\eta-(\alpha+\alpha')\xi}{\dd}
 =\frac{2\alpha\eta+(B/A)\xi}{\dd}.
\]
Substituting $\xi=Q/(A\eta)$ gives the first identity and the fourth; the second is a
restatement, and $H>0$ forces $\sgn G=\sgn(\eta\dd)$. For the third,
$\varepsilon=|\xi|/H=\bigl(|Q|/(A|\eta|)\bigr)\cdot A^2|\eta|\,|\dd|/|G|=A|\dd|\,|Q|/|G|$ and
$A|\dd|=\sqrt\Delta$. For \eqref{eq:Qz}, solve the displayed $H$-identity for $\eta$ and
substitute into $Q=A\xi\eta$ with $\xi=zH$; the bound follows from
$|A\dd-Bz|\le A|\dd|+|B|\,|z|$.
\end{proof}

\begin{definition}[Records]\label{def:records}
A state is a \emph{strict record} if $\varepsilon$ is strictly smaller than at every admissible
state of smaller $H$, and a \emph{weak record} if $\le$ holds. A state whose error exactly
equals the running minimum at a strictly larger height is a weak record that is not strict; we
call such a coincidence a \emph{tie}.
\end{definition}

\begin{remark}[$H$ is a strict total order]\label{rem:order}
If $m\alpha+k=m'\alpha+k'$ with $(m,k)\ne(m',k')$ then $(m-m')\alpha=k'-k$, which forces
$m=m'$ and then $k=k'$ because $\alpha$ is irrational. So distinct states have distinct heights
and $\adm$ is linearly ordered by $H$ with no ties in the ordering; the running minimum of
$\varepsilon$ is therefore well defined, and Definition~\ref{def:records} is unambiguous.
\end{remark}

\subsection{What the objective is, and is not}\label{sec:objective}

Since $\varepsilon=|m\alpha-k|/(m\alpha+k)=|x-\alpha|/(x+\alpha)$ for $x=k/m$, and
$(x-\alpha)/(x+\alpha)=\tanh\bigl(\tfrac12\log(x/\alpha)\bigr)$, we have
\begin{equation}\label{eq:tanh}
\varepsilon=\Bigl|\tanh\Bigl(\tfrac12\log\frac{k}{m\alpha}\Bigr)\Bigr| .
\end{equation}
So records minimise the multiplicative error $|\log(k/m)-\log\alpha|$ under the weighted height
$m\alpha+k$. This is \emph{not} the classical objective, and no transfer in either direction
should be assumed. What does hold is one inclusion.

\begin{proposition}\label{prop:classical}
Call $(m,k)\in\adm$ a \emph{class best approximation} if $|\xi(m,k)|$ is smaller than at every
admissible state of smaller $H$. Every class best approximation is a strict
$\varepsilon$-record, and the inclusion can be strict.
\end{proposition}

\begin{proof}
If $|\xi|$ strictly decreases while $H$ increases then $\varepsilon=|\xi|/H$ strictly
decreases. For strictness of the inclusion: at $\alpha=\varphi$ with $m,k$ odd, $(7,11)$ is a
strict $\varepsilon$-record while $|3\varphi-5|=0.1459\ldots<0.3262\ldots=|7\varphi-11|$ with
$3\le7$, so $(7,11)$ is not a class best approximation.
\end{proof}

\begin{remark}[Not a claim]\label{rem:noteq}
Whether the inclusion is ever an \emph{equality} of the two infinite sets is not settled here.
At $\alpha=\sqrt3$ with $m,k$ odd the two sets agreed at every height tested, but that is an
observation about a finite range and no equality of infinite sets follows from it; see
Section~\ref{sec:limits}.
\end{remark}

\begin{remark}\label{rem:order2}
Proposition~\ref{prop:classical} compares two notions that are \emph{both} ordered by $H$. The
classical and parity-restricted literature \cite{KimLeeLiao2022,KimLeeLiao2025,KushwahaSarma}
orders by the denominator instead, and no bridge between the two orders is established here;
see Section~\ref{sec:limits}. In particular the witness $(7,11)$ in the proof above speaks only
to the $H$-ordered comparison.
\end{remark}

In the geometric picture of Klein and Arnold, best approximations of a real number are the
lattice points on the \emph{sail} --- the boundary of the convex hull of the lattice points in a
specified cone \cite{Karpenkov2013,SpaldingVeselov}. Since an admissible set here is a union of
cosets of $M\Z^2$, the natural reduction to try is: take the cone, take the admissible points
inside it, and read the records off the vertices of their convex hull. That reduction
\emph{fails}, and it fails on an exact finite witness.

\begin{proposition}[A record need not be a vertex of its coset hull]\label{prop:sail}
Let $\alpha=2+\sqrt2$, i.e. $(A,B,C,s)=(1,-4,2,+1)$, with $M=3$ and
$\mathcal R=\{(1,1)\}$. Then $(7,25)$ is a strict record: there are exactly $42$ admissible
states of smaller height, the best of them is $(4,13)$, and
\[
\varepsilon(7,25)=0.02250545\ldots\ <\ \varepsilon(4,13)=0.02464110\ldots
\]
Nevertheless
\begin{equation}\label{eq:sail}
(7,25)=\tfrac35\,(1,4)+\tfrac15\,(31,106)+\tfrac15\,(1,7),
\end{equation}
an exact identity whose three vertices are admissible, lie in the open cone $k>\alpha m$
together with $(7,25)$ itself, and span a nondegenerate triangle $(\det=90)$. All three
barycentric coordinates are strictly positive, so $(7,25)$ lies \emph{strictly inside} that
triangle and is therefore interior to the convex hull of the admissible points of the cone:
neither a hull vertex nor an edge point.
\end{proposition}

\begin{proof}
Each assertion is a finite exact computation, carried out in
\texttt{scripts/review6\_sail.py} ($14/14$ checks): the enumeration of the $42$ lower states
and the two error comparisons in $\Z[\sqrt2]$; the identity \eqref{eq:sail} over $\Q$; the
residues $(1,4),(31,106),(1,7)\equiv(1,1)\bmod3$; the cone memberships
$k-\alpha m=0.585786\ldots,\ 0.159380\ldots,\ 3.585786\ldots$ and $1.100505\ldots$; and
$\det=90$. Strict positivity of all three barycentric coordinates together with
nondegeneracy is exactly interiority.
\end{proof}

\begin{remark}[What Proposition~\ref{prop:sail} does and does not do]\label{rem:sailscope}
It refutes \emph{this} reduction --- records as hull vertices of the admissible points of a
cone --- on a single exact witness, and no amount of further data can undo a finite witness. It
does \emph{not} rule out other geometric descriptions of the record set, and it does not by
itself establish that the classification below is new. The sails of \cite{SpaldingVeselov} are
moreover attached to a specified pair of lines, a precision that matters when comparing their
construction with the present record problem.
\end{remark}

\section{The main theorem}\label{sec:main}

Throughout, orbits are orbits of the cyclic group $\langle P\rangle$ generated by the matrix of
Lemma~\ref{lem:automorph}, acting on column vectors $(m,k)^{\!\top}$.

\begin{definition}[Levels]\label{def:level}
Each orbit is given a representative with $|\eta|\in[1,\rho)$, where $\rho>1$ is the factor by
which $P$ multiplies $\eta$ (Lemma~\ref{lem:automorph}). The \emph{level} of the state obtained
from the representative by $L$ applications of $P$ is $L\in\Z$; equivalently, $L$ is the unique
integer with $|\eta|/\rho^{L}\in[1,\rho)$. Levels are absolute: there is no per-orbit offset.
\end{definition}

\begin{theorem}[Effective classification of strict records]\label{thm:main}
There is an algorithm which, on input $(A,B,C,s,M,\mathcal R)$ as in Section~\ref{sec:setting},
terminates and outputs
\begin{enumerate}\itemsep2pt
\item a matrix $P\in SL_2(\Z)$ with $P\equiv I\pmod M$ which preserves $Q$ exactly, preserves
$\gcd(m,k)$, and multiplies $\eta$ by a fixed $\rho>1$;
\item a state $w\in\adm$ with attained error $\epsw:=\varepsilon(w)<\epsbar$, where
$\epsbar:=\min\bigl(1,\sqrt\Delta/(2|B|)\bigr)$ and the second entry is omitted when $B=0$;
\item integers $B_{\mathrm{cut}}\le B_{\mathrm{comp}}$ and a finite list of orbits
$O_1,\dots,O_N$ which together contain every admissible state with $|Q|\le B_{\mathrm{comp}}$;
\item for each $O_i$ an integer $L_i^{*}$ together with the two exact witnesses
$Y_{L_i^{*}}\le Y_i^{*}<Y_{L_i^{*}+1}$, and a tail-start level $T_i>L_i^{*}$;
\item for each $O_i$ a verdict in $\{\textsf{record},\textsf{tie},\textsf{excluded}\}$ which is
correct at \emph{every} level $L\ge T_i$;
\item the exact strict-record set of the finite region $H\le H(w)$, and the exact verdict of
each state at the finitely many levels $L_i^{*}<L<T_i$.
\end{enumerate}
The strict records of $\adm$ are then exactly
\begin{align*}
&\underbrace{\{\text{records with }H\le H(w)\}}_{\text{(6), finite}}
\ \cup\
\underbrace{\{\text{records at levels }L_i^{*}<L<T_i\}}_{\text{(6), finite}}\\[2pt]
&\qquad\cup\quad
\bigcup_{i\,:\,O_i\ \textsf{record}}\ \bigl\{\text{states of }O_i\text{ at levels }
L\ge T_i\bigr\},
\end{align*}
and the last set is the complete forward $\langle P\rangle$-orbit of the seed
$P^{\,T_i}(\text{rep}_i)$. Every output item is accompanied by exact inequalities that a
checker can re-derive from the definitions.
\end{theorem}

The three unions are disjoint, and together with the analytically excluded parts they exhaust
$\adm$; this is made precise as Proposition~\ref{prop:partition}. The theorem is about
\emph{strict} records. Ties are detected and reported --- item (5) --- and the certified ties
are genuine infinite families of weak records, but we do \emph{not} claim a classification of
all weak records; see Section~\ref{sec:limits}.

\begin{remark}[Effectiveness versus cost]\label{rem:cost}
Every step below terminates for every admissible input, and that is what Theorem~\ref{thm:main}
asserts. It does \emph{not} assert a bound on the running time. The dominant cost is the
enumeration of Section~\ref{sec:enum}, whose size grows with $\rho=\lambda^{\,j}$ where $j$ is
the index of Lemma~\ref{lem:automorph}, and $j$ is bounded only by $|GL_2(\Z/M\Z)|$. The two
are separated deliberately; see Section~\ref{sec:limits}.
\end{remark}

\section{The class-fixing automorph}\label{sec:automorph}

\begin{lemma}\label{lem:automorph}
Let $(T,U)$ be the least solution of $T^2-\Delta U^2=4$ in positive integers. Then
\[
X_0:\ (m,k)\longmapsto
\Bigl(\tfrac{T+BU}{2}\,m+AU\,k,\ -CU\,m+\tfrac{T-BU}{2}\,k\Bigr)
\]
is an integral matrix of determinant $1$ with $Q(X_0v)=Q(v)$ for all $v$, and $X_0$ multiplies
$\eta$ by the unit $\lambda=(T+U\sqrt\Delta)/2$ of the order of discriminant $\Delta$; after
replacing $X_0$ by $X_0^{-1}$ if necessary, $\lambda>1$. Moreover there is a least
$j\le|GL_2(\Z/M\Z)|$ with
\[
P:=X_0^{\,j}\equiv I \pmod M,
\]
and $P$ preserves $Q$, fixes every residue class modulo $M$, preserves $\gcd(m,k)$, and
multiplies $\eta$ by $\rho:=\lambda^{\,j}>1$.
\end{lemma}

\begin{proof}
Integrality: $T^2\equiv\Delta U^2\equiv B^2U^2\pmod4$, so $T\equiv BU\pmod2$ and both displayed
halves are integers. That $X_0$ is the proper automorph of the form $Q$, of determinant $1$, and
that it acts on $\eta$ as multiplication by $\lambda$, is the classical automorph of an
indefinite binary quadratic form \cite{BuellBQF,ConradGenPell}; $\lambda>1$ after the possible
inversion because $\Delta>0$ and $(T,U)$ is positive. Reducing modulo $M$ puts $X_0$ in the
finite group $GL_2(\Z/M\Z)$, so the cyclic subgroup it generates is finite and some power is the
identity there; take the least such $j$. $P\equiv I\pmod M$ fixes every residue class
individually, hence maps $\adm$-residues to themselves; $P\in GL_2(\Z)$ gives
$\gcd(Pv)=\gcd(v)$; and $Q(Pv)=Q(v)$ and $\eta(Pv)=\rho\,\eta(v)$ by iteration.
\end{proof}

\begin{remark}[Why proper automorphs only]\label{rem:proper}
One may instead use a power of the fundamental \emph{unit}, which may have norm $-1$; then $Q$
alternates sign along an orbit and every comparison below acquires a parity-dependent
right-hand side. Restricting to proper automorphs makes $Q$ \emph{constant} on each orbit at
the cost of splitting each former orbit into at most two. Every comparison below then has a
constant right-hand side and no verdict alternates.
\end{remark}

\begin{remark}[No entrywise-positive power when $\alpha'>0$]\label{rem:nopos}
If $\alpha'>0$ then both eigendirections $(1,\alpha)$ and $(1,\alpha')$ of $P$ lie in the open
positive quadrant, and $P^{n}/\rho^{n}$ tends to a multiple of $(1,\alpha)\,(-\alpha',1)$, which
has a negative entry. Hence \emph{no} power of the automorph is entrywise positive and the
positive quadrant is not forward invariant. For example at $\alpha=(3+\sqrt5)/2$ one gets
$P=\binom{-3\ \ 8}{-8\ \ 21}$. Any argument of the form ``once a state is positive all its
successors are'' is therefore unavailable here, and is replaced by Lemma~\ref{lem:window}.
\end{remark}

\section{A comparison family, unconditionally}\label{sec:family}

\begin{lemma}\label{lem:window}
Along an orbit $\sgn\eta$ is constant and $\xi\to0$ in forward direction. Consequently
$m=(\eta-\xi)/\dd\to+\infty$ forward if and only if $\sgn\eta=\sgn\dd$, and then also
$k=\xi+m\alpha\to+\infty$. Thus:
\begin{enumerate}\itemsep2pt
\item an orbit has an infinite forward tail of positive states iff $\sgn(\eta\,\dd)>0$, and then
$\varepsilon\to0$ along it;
\item if $\alpha'<0$, every orbit either has such a tail or has only finitely many positive
states;
\item if $\alpha'>0$, an orbit with $Q<0$ is positive in \emph{both} directions, and backward
$\varepsilon$ increases to the limit $|\dd|/(\alpha+\alpha')>0$.
\end{enumerate}
\end{lemma}

\begin{proof}
$P$ multiplies $\eta$ by $\rho>0$ and $\xi=Q/(A\eta)$ by $\rho^{-1}$, giving the first
sentence and (1), with $\varepsilon=|\xi|/H\to0$ since $|\xi|\to0$ and $H\to\infty$. For (2)
and (3), in the backward direction $\eta\to0$ and $|\xi|\to\infty$, so
$m\approx-\xi/\dd$ and $k=\xi+m\alpha\approx-\xi\alpha'/\dd$; both are positive exactly when
$\sgn\xi=-\sgn\dd$ and $\alpha'>0$, which with $\sgn(\eta\dd)>0$ means $\xi\eta<0$, i.e.
$Q<0$. On such a ray $H\approx-\xi(\alpha+\alpha')/\dd$ and $\varepsilon=|\xi|/H\to
|\dd|/(\alpha+\alpha')$.
\end{proof}

\begin{lemma}[Existence of a comparison family]\label{lem:fam}
Every nonempty $\mathcal R$ admits an admissible state $A_1$ with $\sgn(\eta\dd)>0$, and hence a
family $A_1,A_2,\dots$ --- its forward $P$-orbit, or the union of finitely many such orbits ---
of admissible states with heights increasing to infinity and $\varepsilon\to0$. No prescribed
or minimal defect is needed: $|Q(A_1)|$ is whatever it is, and is computed.
\end{lemma}

\begin{proof}
Pick $(a,b)\in\mathcal R$. The open interval with endpoints $m\alpha'$ and $m\alpha$ --- that is
$(m\alpha',m\alpha)$ when $\alpha>\alpha'$, and $(m\alpha,m\alpha')$ when $\alpha$ is the
smaller root --- has length $m|\dd|\to\infty$, so for $m\equiv a\pmod M$ large enough it
contains $M$ consecutive integers and hence some $k\equiv b\pmod M$. For such $(m,k)$ the value
$k$ lies strictly between $m\alpha'$ and $m\alpha$, so $\xi$ and $\eta$ have opposite signs and
$\sgn\eta=\sgn\dd$. Lemma~\ref{lem:window}(1) does the rest; $Q\ne0$ and is computable.
\end{proof}

\begin{remark}
Stage~3 of the earlier programme asked for an admissible state of \emph{prescribed} or
\emph{minimal} defect and recorded its absence as a gap. Lemma~\ref{lem:fam} shows the theorem
needs only existence. A small $|Q(A_1)|$, and a family built from several orbits, improve the
constant in Section~\ref{sec:cutoff} but are not required.
\end{remark}

\section{A sufficient record-defect cutoff}\label{sec:cutoff}

\begin{theorem}[Cutoff]\label{thm:cutoff}
Let $A,C$ be consecutive members of an admissible comparison family and let $B^{*}$ be a strict
record with $H_A<H_{B^{*}}\le H_C$. Then
\begin{equation}\label{eq:cut}
|Q_{B^{*}}|\ <\ |Q_A|\Bigl(\frac{H_C}{H_A}\Bigr)^{2}\bigl(1+\theta_A\bigr),
\qquad
\theta_A=\frac{BQ_A+|BQ_A|}{2\alpha A^2\eta_A^2}=2\max(u_A,0).
\end{equation}
\end{theorem}

\begin{proof}
$B^{*}$ is a record above $A$, so $\varepsilon_{B^{*}}<\varepsilon_A$; the map
$x\mapsto x(\sqrt\Delta+|B|x)$ is strictly increasing on $x\ge0$; and $H_{B^{*}}\le H_C$.
Hence \eqref{eq:Qz} gives
$|Q_{B^{*}}|<H_C^{2}\varepsilon_A(\sqrt\Delta+|B|\varepsilon_A)/(2\alpha)$. Now substitute
$\varepsilon_A=|Q_A|\sqrt\Delta/G_A$ and $H_C^2=G_C^2/(Y_C\Delta)$ from \eqref{eq:ident} and
simplify, using $Q_C=Q_A$ and $Y_C=\rho^2Y_A$ for a single-orbit family, to reach
\eqref{eq:cut}.
\end{proof}

No sign hypothesis is used: not $B\le0$, not $\alpha'<0$, not $\alpha>1$, not $A=1$, not
primitivity.

\begin{lemma}[A computable majorant]\label{lem:tailbound}
Write $u=BQ/(2\alpha Y)$, so $u$ is multiplied by $\rho^{-2}$ at each step. For a consecutive
family pair with $r=|\eta_C/\eta_A|$ and $v=\max(|u_A|,|u_C|)<1$,
\begin{equation}\label{eq:maj}
|Q_A|\Bigl(\frac{H_C}{H_A}\Bigr)^{2}(1+\theta_A)
\ \le\ |Q_A|\,r^{2}\Bigl(\frac{1+v}{1-v}\Bigr)^{2}(1+2v).
\end{equation}
Consequently
$B_{\mathrm{cut}}:=\bigl\lfloor\max\bigl(\text{\eqref{eq:cut} over the finitely many pairs
before a cut level},\ \text{\eqref{eq:maj} at that level}\bigr)\bigr\rfloor$
is a computable integer, and every strict record of height $>H(A_1)$ satisfies
$|Q|\le B_{\mathrm{cut}}$.
\end{lemma}

\begin{proof}
By \eqref{eq:ident}, $H_C/H_A=(\eta_C/\eta_A)(1+u_C)/(1+u_A)$, so
$(H_C/H_A)^2\le r^2((1+v)/(1-v))^2$, and $\theta_A=2\max(u_A,0)\le 2v$. Since $|u|$ decreases
by the factor $\rho^{-2}$ along the family, \eqref{eq:maj} evaluated at a cut level bounds every
later pair, so the displayed maximum bounds \eqref{eq:cut} over all pairs.
\end{proof}

\begin{warning}\label{warn:sup}
$B_{\mathrm{cut}}$ is a \emph{certified upper bound} obtained from the majorant \eqref{eq:maj};
it is not claimed to equal the supremum of \eqref{eq:cut}. Only the implication
``record $\Rightarrow|Q|\le B_{\mathrm{cut}}$'' is used.
\end{warning}

\begin{remark}[Families beat single ladders]\label{rem:families}
Each family member contributes at least $|Q_A|$ to the maximum, so
$B_{\mathrm{cut}}\ge\max_A|Q_A|$: the cutoff cannot be pushed below the largest defect the
family contains. The useful trade-off is therefore to take the family to be \emph{all}
admissible orbits of defect $\le D$ for a small $D$, obtained from the enumeration of
Section~\ref{sec:enum} with no height search, and to minimise over a short list of $D$. The
effect can be large: at $\alpha=\sqrt6$ with $(m,k)\equiv(1,1)\pmod3$, where
$\rho=(5+2\sqrt6)^6=940898$, a single-orbit ladder gives a cutoff near $7.3\cdot10^{8}$,
whereas the family of all orbits with $|Q|\le50$ gives $B_{\mathrm{cut}}=135$.
\end{remark}

\section{Complete finite enumeration}\label{sec:enum}

\begin{proposition}[Reduction box]\label{prop:box}
Fix $B_{\max}>0$. Normalise $|\eta|$ into $[1,\lambda)$ by powers of $X_0$. Then
$|\xi|=|Q|/(A|\eta|)\le B_{\max}/A$, and since $m=(\eta-\xi)/\dd$ and $k=\xi+m\alpha$,
\[
|m|\le\frac{\lambda+B_{\max}/A}{|\dd|}+1,\qquad
|k|\le\frac{B_{\max}}{A}+|m|\alpha+1 .
\]
An exhaustive scan of that box meets every $\lambda$-class of states with $0<|Q|\le B_{\max}$.
\end{proposition}

This is the classical reduction theory of the generalised Pell equation
\cite{ConradGenPell,BuellBQF,Cohen1993}: finiteness of the classes of solutions of a given norm
under the automorph group, together with an explicit bounded region containing a representative
of each. \emph{No class-number or principal-ideal hypothesis is used}, and representatives with
non-positive coordinates must be retained, since their orbits enter the positive region
elsewhere.

\begin{proposition}[Splitting and admissibility]\label{prop:split}
Each $\lambda$-class splits into at most $j$ $P$-orbits. Walking $p=0,\dots,j-1$ from the class
representative by $X_0$ and keeping those translates whose residue lies in $\mathcal R$
enumerates exactly the admissible $P$-orbits with $0<|Q|\le B_{\max}$.
\end{proposition}

Using $\lambda$ rather than $\rho=\lambda^{\,j}$ for the box is essential: the box size is
$O(\lambda+B_{\max})$ instead of $O(\lambda^{\,j}+B_{\max})$.

\begin{lemma}[Imprimitive states are retained]\label{lem:imprim}
$P\in GL_2(\Z)$ preserves $\gcd$, so the orbit of an imprimitive state is a separate orbit
inside $\adm$ with the same gcd, and $|Q|\le B_{\max}$ bounds $g=\gcd(m,k)$ by
$\sqrt{B_{\max}/|Q_0|}$ where $Q_0$ is the defect of the reduced pair. Imprimitive states are
\emph{not} discardable in general.
\end{lemma}

\begin{warning}\label{warn:imprim}
For the odd/odd class one may discard them, because the gcd of two odd numbers is odd and
dividing it stays in the class. That argument is special. At $\alpha=\sqrt6$ with
$(m,k)\equiv(1,1)\pmod3$ the state $(4,10)$ has $\gcd=2$ and defect $+4$, is a strict record,
and its reduction $(2,5)$ is $(2,2)\bmod3\notin\mathcal R$; it generates a \emph{permanent
record tail}, and three further imprimitive record tails occur in the same instance. At
$\alpha=\sqrt3$ with $\mathcal R=\{(1,1),(2,3)\}\bmod5$ there are record tails of $\gcd$ $2$
and $3$.
\end{warning}

\section{The comparison state and the blocker bound}\label{sec:eps0}

Two symbols are kept distinct: the \emph{bound}
\[
\epsbar:=\min\Bigl(1,\ \frac{\sqrt\Delta}{2|B|}\Bigr)\quad(\text{second entry omitted if }B=0),
\]
and the \emph{attained} value $\epsw=\varepsilon(w)$ of a chosen comparison state $w$ with
$\epsw<\epsbar$. Such a $w$ exists by Lemma~\ref{lem:fam}, since $\varepsilon\to0$ along a
forward tail; the implementation maximises $\epsw$ over an explicit finite candidate family ---
the first state of each family seed's forward tail that falls below $\epsbar$ --- because a
larger $\epsw$ gives a shallower cutoff level in Lemma~\ref{lem:cutlevel}. \emph{No global
maximum over $\adm$ is claimed or needed: any computable admissible $w$ with
$\varepsilon(w)<\epsbar$ serves.} The tail analysis is conducted only above $H(w)$; the region
$H\le H(w)$ is finite and is decided directly.

\begin{lemma}[Small error forces $G>0$]\label{lem:G}
If $\varepsilon(m,k)<\sqrt\Delta/|B|$ then $\eta\,\dd>0$, hence $G>0$.
\end{lemma}

\begin{proof}
By Lemma~\ref{lem:ident}, $\eta=(H\dd-(B/A)\xi)/(2\alpha)$ and $|\xi|=\varepsilon H$, so
$\eta\,\dd\ge H|\dd|(|\dd|-|B|\varepsilon/A)/(2\alpha)$, which is positive exactly when
$\varepsilon<A|\dd|/|B|=\sqrt\Delta/|B|$. Then $G=A^2\eta\,\dd\,H>0$.
\end{proof}

Since $\epsw<\epsbar\le\sqrt\Delta/(2|B|)$, every strict record above $H(w)$ and every state
that blocks one has error below $\epsw$ and therefore lies on the branch $G>0$. This is what
makes the error monotone (Lemma~\ref{lem:epsmono}) and the height eventually increasing
(Section~\ref{sec:compare}) wherever those facts are used.

\begin{lemma}[Blocker bound]\label{lem:blocker}
Say $V$ \emph{blocks} the state $B^{*}$ if $H_V<H_{B^{*}}$ and
$\varepsilon_V\le\varepsilon_{B^{*}}$. If $V$ blocks $B^{*}$ and
$\sqrt\Delta>|B|\varepsilon_{B^{*}}$ then
\[
|Q_V|\ \le\ |Q_{B^{*}}|\,
\frac{\sqrt\Delta+|B|\varepsilon_{B^{*}}}{\sqrt\Delta-|B|\varepsilon_{B^{*}}} .
\]
In particular, when $\varepsilon_{B^{*}}\le\sqrt\Delta/(2|B|)$ --- which holds for every record
above $H(w)$ --- the bound is $|Q_V|\le3|Q_{B^{*}}|$, so
$B_{\mathrm{comp}}:=3B_{\mathrm{cut}}$ is always safe.
\end{lemma}

\begin{proof}
Apply \eqref{eq:Qz} to $V$ and to $B^{*}$: $|Q_V|\le H_V^2\varepsilon_V(\sqrt\Delta+|B|
\varepsilon_V)/(2\alpha)$, which is at most
$H_{B^{*}}^2\varepsilon_{B^{*}}(\sqrt\Delta+|B|\varepsilon_{B^{*}})/(2\alpha)$; and
$|Q_{B^{*}}|\ge H_{B^{*}}^2\varepsilon_{B^{*}}(\sqrt\Delta-|B|\varepsilon_{B^{*}})/(2\alpha)$
from the same identity. Divide.
\end{proof}

The gain from $\epsbar$ is that the bound of Lemma~\ref{lem:blocker} is \emph{never vacuous};
the implementation uses the tighter per-candidate value and records both. The uniform
$3B_{\mathrm{cut}}$ is what the proof needs.

\section{The two backward exclusions}\label{sec:backward}

\begin{lemma}[Error monotonicity, with its branch]\label{lem:epsmono}
On the branch $G>0$, $\varepsilon$ strictly decreases along an orbit as $|\eta|$ grows. On the
branch $G<0$ it strictly \emph{increases}, and there $\varepsilon\ge\sqrt\Delta/|B|$.
\end{lemma}

\begin{proof}
$\varepsilon=|Q|\sqrt\Delta/|G|$ with $G=2\alpha Y+BQ$ strictly increasing in $Y=A^2\eta^2$.
Where $G>0$, $|G|$ grows and $\varepsilon$ falls. Where $G<0$, $|G|=|BQ|-2\alpha Y$ falls, so
$\varepsilon$ grows, and $|G|\le|BQ|$ gives
$\varepsilon\ge|Q|\sqrt\Delta/|BQ|=\sqrt\Delta/|B|$.
\end{proof}

\begin{warning}\label{warn:epsmono}
The unconditional statement ``$\varepsilon$ decreases along every orbit'' is false. At
$\alpha=(3+\sqrt5)/2$ with $m,k$ odd and $P=\binom{-3\ \ 8}{-8\ \ 21}$ we have
$P(55,21)=(3,1)$, both states positive with $Q=1$, yet
\[
\varepsilon(55,21)=0.745442\ldots<\varepsilon(3,1)=0.774116\ldots,
\]
because $G=-2.9997$ and $-2.8885$ there. Both respect the floor $\sqrt5/3=0.745356\ldots$ of
Lemma~\ref{lem:epsmono}.
\end{warning}

\begin{lemma}[Backward exclusion, two cases]\label{lem:backward}
No state in either case below is a strict record of height $>H(w)$, nor blocks one.
\begin{enumerate}\itemsep2pt
\item \textbf{Orbits with $\sgn(\eta\dd)<0$.} There $G<0$ at every positive state, so
$\varepsilon\ge\sqrt\Delta/|B|\ge2\epsbar>\epsw$ by Lemma~\ref{lem:epsmono}.
\item \textbf{Doubly-infinite orbits ($G>0$ throughout).} Here the floor of case~(1) is
\emph{false}. Instead, on $G>0$,
\begin{equation}\label{eq:ystar}
\varepsilon\ \ge\ \epsw
\quad\Longleftrightarrow\quad
Y\ \le\ Y^{*}:=\frac{|Q|\sqrt\Delta/\epsw-BQ}{2\alpha},
\end{equation}
an explicit cutoff on the level. Levels with $Y\le Y^{*}$ are excluded; only the finitely many
levels with $Y>Y^{*}$ below the tail start remain.
\end{enumerate}
\end{lemma}

\begin{proof}
Case (1) is Lemma~\ref{lem:epsmono} together with
$\sqrt\Delta/|B|\ge2\epsbar>\epsw$. In case (2), $G>0$ gives
$\varepsilon=|Q|\sqrt\Delta/(2\alpha Y+BQ)$, and $\varepsilon\ge\epsw$ rearranges to
\eqref{eq:ystar}. In both cases a state with $H>H(w)$ and $\varepsilon\ge\epsw=\varepsilon(w)$
is not a strict record, because $w$ has smaller height and error at most its own; and it cannot
block a record $B^{*}$ above $H(w)$, since such a $B^{*}$ has
$\varepsilon_{B^{*}}<\epsw\le\varepsilon$.
\end{proof}

\begin{warning}\label{warn:floor}
The case split is necessary: the floor of case (1) does not cover case (2). At
$\alpha=(3+\sqrt5)/2$ with $m,k$ odd, $P^{-1}(1,1)=(13,5)$ lies on a doubly-infinite positive
orbit with $Q=-1$ and $G=3.006211>0$, yet
\[
\varepsilon(13,5)=0.743816\ldots\ <\ \frac{\sqrt\Delta}{|B|}=\frac{\sqrt5}{3}=0.745356\ldots,
\]
and every further backward state of that ray is also below the floor: on the backward ray
$\varepsilon$ rises to $\sqrt\Delta/|B|$ \emph{from below}.
\end{warning}

\begin{lemma}[The cutoff level, with termination]\label{lem:cutlevel}
Fix an orbit with normalised $|\eta|=r\in[1,\rho)$ and put $Y_L=A^2r^2\rho^{2L}$, the value of
$Y$ at level $L$. Since $\rho>1$, $Y_L$ is strictly increasing in $L$ with $Y_L\to0$ as
$L\to-\infty$ and $Y_L\to\infty$ as $L\to+\infty$. As $Y^{*}$ is finite and positive there is a
\emph{unique} integer $L^{*}$ with
\begin{equation}\label{eq:witness}
Y_{L^{*}}\ \le\ Y^{*}\ <\ Y_{L^{*}+1},
\end{equation}
and a two-sided expanding exact search locates it in finitely many steps. The two inequalities
\eqref{eq:witness} are the whole certificate: by monotonicity they prove $Y_L\le Y^{*}$, hence
$\varepsilon\ge\epsw$, for \emph{every} $L\le L^{*}$, and $Y_L>Y^{*}$ for every $L>L^{*}$. The
orbit list is finite, so all such searches terminate.
\end{lemma}

Written logarithmically, $L^{*}=\bigl\lfloor\log(Y^{*}/(A^2r^2))/(2\log\rho)\bigr\rfloor$; the
implementation and the certificate use exact power comparisons only, so no rounding enters.

\begin{lemma}[The retained region is at or beyond the height minimum]\label{lem:ymin}
If $BQ>0$ then $H$, as a function of $Y$ along the orbit, has its minimum at
$Y_{\min}=BQ/(2\alpha)$, and $Y^{*}\ge Y_{\min}$. Hence on the retained region $Y>Y^{*}$ one
has $Y>Y_{\min}$ and the height is strictly increasing in the level.
\end{lemma}

\begin{proof}
$H=2\alpha\eta/\dd+BQ/(A^2\eta\dd)$ has $\partial H/\partial\eta=(2\alpha A^2\eta^2-BQ)/
(A^2\eta^2\dd)$, vanishing at $Y=BQ/(2\alpha)$. And $Y^{*}\ge Y_{\min}$ reduces to
$|Q|\sqrt\Delta/\epsw\ge2BQ$, which holds since $\epsw\le\sqrt\Delta/(2|B|)$ and
$BQ\le|B||Q|$. On $Y>Y^{*}\ge Y_{\min}$ the derivative is positive.
\end{proof}

\begin{remark}
Using the attained $\epsw$ rather than the bound $\epsbar$ gives extra margin and a shallower
$L^{*}$, which is why the implementation does it. It is not needed for
Lemma~\ref{lem:ymin}: the retained condition is the strict $Y>Y^{*}$, so even in the boundary
case $Y^{*}=Y_{\min}$ the conclusion holds.
\end{remark}

\section{Exact comparison formulas}\label{sec:compare}

Because $P$ is a proper automorph, $Q$ is \emph{constant} along each orbit. Write $r_A\in[1,\rho)$
for the normalised $|\eta|$ of orbit $A$.

\begin{lemma}[Error comparison]\label{lem:epscmp}
Let $G_A,G_B>0$, with $A$ at level $p-d$ and $B$ at level $p$. Then
\begin{equation}\label{eq:crit}
\varepsilon_B<\varepsilon_A
\iff |Q_A|G_B-|Q_B|G_A>0
\iff 2\alpha A^2\Gamma>B\bigl(|Q_B|Q_A-|Q_A|Q_B\bigr),
\end{equation}
where $\Gamma=|Q_A|\eta_B^2-|Q_B|\eta_A^2=\rho^{2(p-d)}\Gamma_d$ and
$\Gamma_d=|Q_A|r_B^2\rho^{2d}-|Q_B|r_A^2$ is independent of $p$. The right-hand side is a
constant integer bounded in absolute value by $2|B||Q_A||Q_B|$, so if $\Gamma_d\ne0$ then
$\sgn\Gamma_d$ decides \eqref{eq:crit} for all $p$ with
\begin{equation}\label{eq:epsthr}
2\alpha A^2\rho^{2(p-d)}|\Gamma_d|>\bigl|B(|Q_B|Q_A-|Q_A|Q_B)\bigr|,
\end{equation}
an explicitly computable threshold. If $\Gamma_d=0$ the verdict is the constant sign of
$-B(|Q_B|Q_A-|Q_A|Q_B)$: a three-way outcome, and for $B=0$ an \emph{exact permanent equality}.
\end{lemma}

\begin{proof}
$\varepsilon=|Q|\sqrt\Delta/G$ on $G>0$ gives the first equivalence; expanding
$G=2\alpha A^2\eta^2+BQ$ gives $2\alpha A^2\Gamma+B(|Q_A|Q_B-|Q_B|Q_A)>0$, which is the second.
$\Gamma$ scales by $\rho^2$ per level because $Q$ is constant, whence the factorisation.
\end{proof}

\begin{warning}\label{warn:sign}
The right-hand side of \eqref{eq:crit} is $B(|Q_B|Q_A-|Q_A|Q_B)$ and not its negative. A single
exact instance pins the orientation and is kept as a regression: at $\alpha=\varphi$ with
$m,k$ odd take $a=(1,1)$, $b=(1,3)$, so $Q_a=-1$, $Q_b=5$ and $\Gamma=0$ exactly, with
$G_a,G_b>0$. The correct right-hand side is $+10$ and $0>10$ is false, predicting
$\varepsilon_b\ge\varepsilon_a$; indeed
\[
\varepsilon_a=\sqrt5-2=0.236068\ldots<\varepsilon_b=\tfrac{10-3\sqrt5}{11}=0.299254\ldots
\]
The opposite sign would read $0>-10$ and predict the reverse.
\end{warning}

\begin{lemma}[Height comparison]\label{lem:hcmp}
Let $A$ be at level $q$ and $B$ at level $q+d$, on the branch $\eta\dd>0$. Then
\begin{equation}\label{eq:hcrit}
\sgn(H_B-H_A)=\sgn\bigl(\rho^{2q}U+V\bigr),\qquad
\begin{aligned}
U&=2\alpha A^2r_Ar_B\rho^{d}\bigl(r_B\rho^{d}-r_A\bigr),\\
V&=B\bigl(Q_Br_A-Q_Ar_B\rho^{d}\bigr),
\end{aligned}
\end{equation}
certified once $\rho^{2q}|U|>|V|$. Moreover $U=0$ forces $r_A=r_B\rho^{d}$, i.e. the same orbit
after an index shift, and then $V=0$ as well.
\end{lemma}

\begin{proof}
$H=2\alpha\eta/\dd+BQ/(A^2\eta\dd)$. On the branch $\eta\dd>0$ the quantity
$A^2|\dd|\,r_Ar_B\rho^{q+d}$ is \emph{positive}; multiplying $H_B-H_A$ by it turns the first
term into $\rho^{2q}U$ and the second into exactly $V$. If $r_A=r_B\rho^d$ then taking norms
gives $Q_A=Q_B$ and $V=B(Q_Br_A-Q_Ar_B\rho^d)=BQ_B(r_A-r_B\rho^d)=0$.
\end{proof}

\begin{warning}\label{warn:hsign}
Two details of \eqref{eq:hcrit} are easy to get wrong. The exponent is $\rho^{2q}$, not
$\rho^{q}$: at $\alpha=2+\sqrt2$ with $X_0=\binom{-1\ \ 2}{-4\ \ 7}$, $\rho=3+2\sqrt2$ and
normalised representatives $a_0=(-61,-30)$, $b_0=(-59,-29)$ mapping to $a_1=(1,34)$,
$b_1=(1,33)$, one has $H_{b_1}-H_{a_1}=-1$ exactly, while at $q=1,d=0$
\[
\rho^{q}U+V=-7196+5460\sqrt2=+525.606\ldots\ (\text{wrong sign}),\qquad
\rho^{2q}U+V=-525.606\ldots\ (\text{right}).
\]
And the clearing factor must carry $|\dd|$: for the smaller root $\dd<0$ --- e.g.
$\dd=-2.236068$ at $\alpha=(3-\sqrt5)/2$ --- while $\eta\dd>0$ on the retained states.
\end{warning}

\section{A finite comparison ledger}\label{sec:ledger}

For each ordered pair of orbits and each offset $d$, Lemmas~\ref{lem:epscmp} and \ref{lem:hcmp}
supply one exact sign and one threshold. To make the collection finite both endpoints of the
offset range must be justified.

\begin{lemma}[Upper endpoint]\label{lem:dstar}
Let $d^{*}$ be the least $d\ge1$ with $\rho^{2d-2}>B_{\mathrm{comp}}$. On the branch $G>0$,
$\varepsilon_A>\varepsilon_{B^{*}}$ holds as soon as $|Q_A|G_{B^{*}}>|Q_{B^{*}}|G_A$, and using
$r_A<\rho$, $r_{B^{*}}\ge1$ it suffices that
\begin{equation}\label{eq:noblock}
2\alpha A^2\rho^{2L}\bigl(|Q_A|-|Q_{B^{*}}|\rho^{2-2d}\bigr)
>|Q_A||BQ_{B^{*}}|+|Q_{B^{*}}||BQ_A| .
\end{equation}
The bracket is positive for every pair with $|Q|\le B_{\mathrm{comp}}$ once $d>d^{*}$, and the
additive part then fixes a level threshold $p_{\mathrm{noblock}}$. Hence for $d>d^{*}$ no
competitor can block at any level $\ge p_{\mathrm{noblock}}$, irrespective of the height
relation, and only $d\le d^{*}$ needs the thresholds \eqref{eq:epsthr}.
\end{lemma}

\begin{lemma}[Lower endpoint]\label{lem:lower}
For every ordered pair the $d=1$ instance of \eqref{eq:hcrit} with the \emph{candidate} in the
first slot certifies $H_{\text{comp}}(q+1)>H_{\text{cand}}(q)$ beyond its threshold. With
Lemma~\ref{lem:ymin} --- the height is strictly increasing in the level on the retained region
--- every competitor state at a level above the candidate's lies above it in height and
therefore cannot block. So negative offsets are excluded and the ledger covers exactly
$0\le d\le d^{*}$.
\end{lemma}

Accordingly put
\begin{align*}
p_{\mathrm{cert}}:=\max\bigl(\,&\text{thresholds of Lemmas~\ref{lem:epscmp} and
  \ref{lem:hcmp}, over every ordered pair with a candidate}\\
&\text{side and every }0\le d\le d^{*};\quad
p_{\mathrm{noblock}};\quad L_{\epsw}\,\bigr),
\end{align*}
where $L_{\epsw}$ is the least level at which every candidate has $\varepsilon<\epsw$. Beyond
$p_{\mathrm{cert}}$ every sign in the ledger is constant.

\section{Binding-state sufficiency}\label{sec:binding}

\begin{lemma}[Sufficiency]\label{lem:binding}
Fix a candidate $B^{*}$ and an orbit $O$. Among the states of $O$ of height below $H_{B^{*}}$,
let $V^{*}$ be one of minimal $\varepsilon$. If \emph{any} state of $O$ blocks $B^{*}$ then
$V^{*}$ blocks $B^{*}$. Moreover the three-way verdict read off $V^{*}$ equals the verdict of an
exhaustive scan of $O$: if $\varepsilon_{V^{*}}>\varepsilon_{B^{*}}$ no state of $O$ blocks; if
$\varepsilon_{V^{*}}=\varepsilon_{B^{*}}$ every blocker in $O$ ties; if
$\varepsilon_{V^{*}}<\varepsilon_{B^{*}}$ there is a strict loss. Hence testing $V^{*}$ alone is
sufficient.
\end{lemma}

\begin{proof}
If $V\in O$ blocks then $\varepsilon_{V^{*}}\le\varepsilon_V\le\varepsilon_{B^{*}}$ and
$H_{V^{*}}<H_{B^{*}}$, so $V^{*}$ blocks. If $\varepsilon_{V^{*}}=\varepsilon_{B^{*}}$ then any
blocker $V$ has $\varepsilon_{B^{*}}\ge\varepsilon_V\ge\varepsilon_{V^{*}}=\varepsilon_{B^{*}}$,
so $\varepsilon_V=\varepsilon_{B^{*}}$ and all blockers tie. The remaining two cases are
immediate.
\end{proof}

\begin{warning}\label{warn:binding}
The stronger-sounding statement ``a state below the $\varepsilon$-minimal one cannot block'' is
\emph{false}, and Lemma~\ref{lem:binding} is deliberately a sufficiency claim. At
$\alpha=\varphi$ with no congruence restriction take $b=(15,24)=3\cdot(5,8)$ and the orbit
$(5,8)\mapsto(13,21)\mapsto(34,55)$ of $\binom{1\ \ 1}{1\ \ 2}$. Then
\[
H(5,8)<H(13,21)<H(b)<H(34,55),\qquad
\varepsilon(13,21)<\varepsilon(5,8)=\varepsilon(b),
\]
so $(13,21)$ is the $\varepsilon$-minimal state below $H(b)$ --- yet the \emph{earlier} $(5,8)$
also blocks $b$, by exact equality of the error. On the branch $G>0$ the
$\varepsilon$-minimal state is the one of largest level (Lemma~\ref{lem:epsmono}), which is why
a largest-level rule also works there; on the branch $G<0$ that rule is reversed, so the
implementation selects by $\varepsilon$ and needs no monotonicity hypothesis.
\end{warning}

\section{The partition, and termination}\label{sec:partition}

Levels are absolute (Definition~\ref{def:level}). For orbit $i$ let $\mathrm{enter}_i$ be the
least non-negative level at which the state is positive, and define the \emph{tail-start level}
\[
T_i:=\mathrm{enter}_i+p_{\mathrm{cert}},\qquad\text{certified to satisfy}\qquad T_i>L_i^{*}.
\]

\begin{proposition}[Exhaustive disjoint partition]\label{prop:partition}
Every admissible state lies in exactly one of the following, and each split is a dichotomy on an
explicitly computable quantity.
\begin{center}\small
\begin{tabular}{@{}llp{7.4cm}@{}}
\toprule
part & condition & status \\
\midrule
$P_1$ & $H\le H(w)$ & decided by direct scan of an explicitly bounded region \\
$P_{2a}$ & $H>H(w)$, $|Q|>B_{\mathrm{comp}}$ & neither a strict record nor a blocker of one
  (Theorem~\ref{thm:cutoff}, Lemma~\ref{lem:blocker}) \\
$P_{2b}$ & $H>H(w)$, $B_{\mathrm{cut}}<|Q|\le B_{\mathrm{comp}}$ & not a record; retained as a
  competitor \\
$P_3$ & $H>H(w)$, $|Q|\le B_{\mathrm{cut}}$, no forward tail & excluded by
  Lemma~\ref{lem:backward}(1) \\
$P_4$ & $H>H(w)$, $|Q|\le B_{\mathrm{cut}}$, forward tail, $L\le L_i^{*}$ & excluded by
  Lemma~\ref{lem:backward}(2) \\
$P_5$ & $H>H(w)$, $|Q|\le B_{\mathrm{cut}}$, forward tail, $L_i^{*}<L<T_i$ & finitely many
  levels (possibly negative $L$), decided exactly \\
$P_6$ & $H>H(w)$, $|Q|\le B_{\mathrm{cut}}$, forward tail, $L\ge T_i$ & the certified tail: one
  verdict per orbit \\
\bottomrule
\end{tabular}
\end{center}
\end{proposition}

\begin{proof}
$P_1$ versus the rest is a height dichotomy. Inside $H>H(w)$, the three-way split on the
integer $|Q|$ against $B_{\mathrm{cut}}\le B_{\mathrm{comp}}$ is exhaustive and disjoint. Inside
$|Q|\le B_{\mathrm{cut}}$, ``forward tail or not'' is decided by $\sgn(\eta\dd)$
(Lemma~\ref{lem:window}). Inside the forward-tail case, $L\le L_i^{*}$ or not is decided by the
two witnesses \eqref{eq:witness}, and then $L<T_i$ or not; since $T_i>L_i^{*}$ is certified,
$P_5$ is well defined (possibly empty) and $P_4,P_5,P_6$ are disjoint. The stated status of each
part is the cited lemma.
\end{proof}

The parts that can be \emph{infinite} are exactly $P_{2a}$, $P_{2b}$, $P_3$ and $P_4$ --- a
doubly-infinite orbit has infinitely many positive states --- and each is excluded by one
inequality, not enumerated. What remains finite is $P_1$, bounded by
\begin{equation}\label{eq:P1box}
1\le m\le\bigl\lfloor H(w)/\alpha\bigr\rfloor,\qquad
1\le k\le\bigl\lfloor H(w)\bigr\rfloor
\end{equation}
followed by the exact height and residue filters, and $P_5$, finitely many levels on finitely
many orbits. $P_6$ is infinite but carries one verdict per orbit.

\begin{remark}[How $P_1$ and $P_5$ completeness is verified]\label{rem:completeness}
Correctness of a \emph{listed} exception is weaker than completeness of the \emph{list}. Both
regions are therefore enumerated independently and compared by \emph{set equality}: $P_1$ from
the explicit bounds \eqref{eq:P1box}, with its record set recomputed as a running minimum; and
$P_5$ by walking the levels $L_i^{*}<L<T_i$ of each candidate orbit and recomputing each verdict
with no pruning. Demanding this found a genuine classification error in an earlier
implementation, which had been reporting $P_1$ and $P_4$ states as $P_5$ anomalies.
\end{remark}

\begin{proposition}[Termination]\label{prop:term}
Every search in Theorem~\ref{thm:main} terminates: the index $j$ of
Lemma~\ref{lem:automorph} is at most $|GL_2(\Z/M\Z)|$; the cutoff of
Lemma~\ref{lem:tailbound} is a maximum over finitely many pairs and one majorant; the box of
Proposition~\ref{prop:box} is finite; the cutoff-level search of Lemma~\ref{lem:cutlevel}
terminates by monotonicity of $Y_L$ and finiteness of $Y^{*}$; the thresholds of
Lemmas~\ref{lem:epscmp}, \ref{lem:hcmp}, \ref{lem:dstar} are each reached after finitely many
exactly tested levels; and the finite regions $P_1$, $P_5$ are bounded by \eqref{eq:P1box} and
by $T_i-L_i^{*}-1$.
\end{proposition}

\begin{remark}[Exact ordering, with a guaranteed fallback]\label{rem:sort}
Every height-ordered list is produced by presorting with an approximate key, verifying
\emph{all} adjacent pairs exactly, and --- if key generation raises or any adjacent pair fails
--- re-sorting with the exact algebraic comparator. Because $H$ is a strict total order
(Remark~\ref{rem:order}), adjacent-pair correctness implies the whole list is sorted, so the
verification is a complete certificate of the ordering. Because of the fallback, verification
failure or key degeneracy costs a slower sort rather than an abort, so termination with a
correctly ordered list holds for every input; verification alone would only prevent a wrong
answer. No ordering decision is ever taken from an approximate key.
\end{remark}

\section{A worked example}\label{sec:example}

Take $\alpha=\varphi=(1+\sqrt5)/2$, i.e. $(A,B,C,s)=(1,-1,-1,+1)$, with $M=2$ and
$\mathcal R=\{(1,1)\}$: both coordinates odd. The algorithm produces the following, all of it
present in the emitted certificate \texttt{control\_phi.json}.

\begin{center}\small
\begin{tabular}{@{}ll@{}}
\toprule
Pell solution & $(T,U)=(3,1)$, $\lambda=\varphi^2=(3+\sqrt5)/2$ \\
fundamental automorph & $X_0=\binom{1\ \ 1}{1\ \ 2}$ \\
class-fixing index & $j=3$, $P=\binom{5\ \ 8}{8\ \ 13}=\varphi^6$, $\rho=17.944272\ldots$ \\
comparison family & all admissible orbits with $|Q|\le1$; lowest member $(1,1)$, $H=2.618\ldots$ \\
cutoff & $B_{\mathrm{cut}}=18$, attained at the pair $\bigl((3,5),(13,21)\bigr)$ \\
comparison state & $w=(1,1)$, $\epsw=\sqrt5-2=0.236068\ldots<\epsbar=1$ \\
blocker bound & $B_{\mathrm{comp}}=22$ (per-candidate; the uniform bound is $3\cdot18=54$) \\
reduction box & $|m|\le13$, $|k|\le45$; $40$ $\lambda$-classes; $14$ admissible orbits \\
certified levels & $d^{*}=2$, $p_{\mathrm{noblock}}=1$, $p_{\mathrm{cert}}=2$ \\
\bottomrule
\end{tabular}
\end{center}

Three orbits receive the verdict \textsf{record}, with first positive states and defects
\[
(1,1)\ [Q=-1,\ L^{*}=0],\qquad
(3,5)\ [Q=+1,\ L^{*}=-1],\qquad
(7,11)\ [Q=-5,\ L^{*}=-1],
\]
each with $T_i=2$. Merging the three tails in increasing height reproduces the sequence
\[
(1,1),\ (3,5),\ (7,11),\ (13,21),\ (55,89),\ (123,199),\ (233,377),\ \dots
\]
whose signed defects cycle $(-1,+1,-5)$; the first family is $(F_{3t+1},F_{3t+2})$ and the third
is the odd Lucas pairs $(L_{6t+4},L_{6t+5})$. One orbit receives the verdict \textsf{tie}: the
imprimitive orbit through $(3,3)=3\cdot(1,1)$, defect $-9$, whose error equals that of the
$(1,1)$ orbit at three times the height, forever. The finite region $H\le H(w)=2.618\ldots$
contains the single record $(1,1)$, and the list of $P_5$ exceptions is empty.

A reader can check one certificate comparison by hand. For the pair $A=(1,1)$, $B=(1,3)$ at
offset $d=0$ one has $Q_A=-1$, $Q_B=5$ and, since $\eta_A=\varphi$, $\eta_B=\varphi+2$,
\[
\Gamma=|Q_A|\eta_B^2-|Q_B|\eta_A^2=(\varphi+2)^2-5\varphi^2=0 ,
\]
so Lemma~\ref{lem:epscmp} falls to its degenerate case: the verdict is the constant sign of
$-B(|Q_B|Q_A-|Q_A|Q_B)=-10$, which is negative, so $(1,3)$ does \emph{not} beat $(1,1)$ --- in
agreement with $\varepsilon(1,1)=0.236068\ldots<0.299254\ldots=\varepsilon(1,3)$. This is
Warning~\ref{warn:sign}.

\section{Stress examples}\label{sec:stress}

The instances below exercise each hypothesis that was removed. All data are read from the
emitted certificates; $\rho=\lambda^{\,j}$.

\begin{center}\footnotesize
\setlength{\tabcolsep}{4pt}
\begin{tabular}{@{}lrrrrrrp{3.5cm}@{}}
\toprule
instance & $j$ & $\rho$ & $B_{\mathrm{cut}}$ & orb. & tails & ties & feature \\
\midrule
$\varphi$, odd/odd & 3 & $17.94$ & 18 & 14 & 3 & 1 & control \cite{DeJesus2026oddodd} \\
$1+\sqrt3$, odd/odd & 2 & $13.93$ & 44 & 31 & 3 & 1 & control \\
$\sqrt3$, odd/odd & 2 & $13.93$ & 27 & 14 & 2 & 4 & control; $B=0$, permanent ties \\
$(3+\sqrt5)/2$, odd/odd & 3 & $17.94$ & 18 & 21 & 3 & 1 & \textbf{two positive conjugates} \\
$(3-\sqrt5)/2$, odd/odd & 3 & $17.94$ & 18 & 21 & 3 & 1 & \textbf{smaller root}, $\alpha<1$ \\
$(5+\sqrt{21})/2$, odd/odd & 3 & $110.0$ & 77 & 102 & 6 & 2 & positive conjugate \\
$1/\sqrt2$, odd/odd & 1 & $5.83$ & 18 & 6 & 2 & 1 & \textbf{nonmonic} $A=2$, $\alpha<1$ \\
$3x^2\!-\!x\!-\!1$, odd/odd & 3 & $1298$ & 32 & 40 & 9 & 3 & \textbf{nonmonic} $A=3$ \\
$\sqrt3$, $\{(1,1),(2,3)\}\bmod5$ & 3 & $51.98$ & 164 & 32 & 6 & 0 & \textbf{residue union};
  two imprimitive \\
$\sqrt6$, $(2,2)\bmod4$ & 2 & $97.99$ & 113 & 12 & 6 & 0 & every state imprimitive;
  all six tails so \\
$\sqrt2$, $(1,1)\bmod3$ & 4 & $1154$ & 61 & 33 & 9 & 2 & $M=3$; two imprimitive tails \\
$\sqrt6$, $(1,1)\bmod3$ & 6 & $9.4\cdot10^{5}$ & 135 & 84 & 20 & 6 & \textbf{imprimitive
  record} $(4,10)$ \\
$N+\sqrt2$, $N=2,\dots,20$ & 1--4 & $5.8$--$1154$ & 15--92 & 17--108 & 2--9 & 0--2 &
  $\sqrt\Delta/|B|=\sqrt2/N\to0$ \\
\bottomrule
\end{tabular}
\end{center}

Three observations. At $\sqrt6$ with $(m,k)\equiv(1,1)\bmod3$ the record tails include four
imprimitive families, among them the seed $(4,10)$ of Warning~\ref{warn:imprim}. The family
$\alpha_N=N+\sqrt2$ has $\Delta=8$ and $\lambda=3+2\sqrt2$ independent of $N$, while the
blocker-bound budget $\sqrt\Delta/|B|=\sqrt2/N$ tends to $0$; it is the regime in which
Lemma~\ref{lem:blocker} is tightest, and $d^{*}\in\{2,3\}$ and $p_{\mathrm{cert}}\in\{2,3\}$
throughout. Across all $26$ instances the cutoff level satisfies $L^{*}\in\{-1,0,1\}$, the
realised blocking offsets are $\{0,1\}$, and the $P_5$ exception lists are all empty.

\section{Reproducibility and verification scope}\label{sec:repro}

\paragraph{What the computations are for.}
Theorem~\ref{thm:main} and every lemma above are proved for arbitrary admissible input. The
computations do three separable things: they instantiate the algorithm on $26$ inputs; they emit
a machine-readable certificate per input containing every quantity the proof requires; and they
audit those certificates with a checker written separately from the classifier. They do
\emph{not} establish the theorem, and no universal statement is inferred from them.

\paragraph{Artifacts.}
The frozen archive accompanying this note contains, with \texttt{sha256} for every file:
\begin{center}\small
\begin{tabular}{@{}ll@{}}
\toprule
\texttt{scripts/gq.py} & exact arithmetic over $\Q(\sqrt\Delta)$ and the algorithm's stages \\
\texttt{scripts/emit.py} & runs the algorithm and writes one certificate per instance \\
\texttt{scripts/check.py} & the independent checker (own arithmetic, no classifier routines) \\
\texttt{scripts/qr\_general.py} & the classifier-side verifier (98 checks) \\
\texttt{scripts/lit\_bridge.py} & the two bridges of Section~\ref{sec:objective} (14 checks) \\
\texttt{scripts/review3,4,5.py} & the regression controls of Warnings~\ref{warn:sign},
  \ref{warn:hsign}, \ref{warn:epsmono}, \ref{warn:floor}, \ref{warn:binding} \\
\texttt{scripts/review6\_sail.py} & the exact witness of Proposition~\ref{prop:sail}
  (14 checks) \\
\texttt{certificates/*.json} & the $26$ certificates \\
\texttt{results/*.log} & the complete logs of the runs quoted below \\
\texttt{notes/*.md} & the working proof and audit notes this note consolidates \\
\bottomrule
\end{tabular}
\end{center}

\paragraph{Commands.} With Python 3.11 or later and no third-party packages:
\begin{center}\small
\begin{tabular}{@{}ll@{}}
\toprule
\texttt{python3 scripts/emit.py} & regenerate all $26$ certificates into \texttt{certificates/} \\
\texttt{python3 scripts/check.py} & audit every certificate in \texttt{certificates/} \\
\texttt{python3 scripts/check.py NAME} & audit one certificate \\
\texttt{python3 scripts/qr\_general.py} & the classifier-side suite \\
\texttt{python3 scripts/review3.py} (4, 5) & the regression controls \\
\texttt{python3 scripts/review6\_sail.py} & the counterexample of
  Proposition~\ref{prop:sail} \\
\texttt{python3 scripts/lit\_bridge.py} & the objective-comparison checks \\
\bottomrule
\end{tabular}
\end{center}

\paragraph{What the checker verifies.}
\texttt{check.py} shares no code, no data representation and no decision routine with the
classifier: it represents elements of $\Q(\sqrt\Delta)$ as $(a+b\sqrt\Delta)/d$ with a single
integer denominator over $\Delta$ itself rather than over its squarefree part, and it implements
its own sign, floor, matrix and comparison routines. It does \emph{not} use the blocker bound:
all competitor comparisons are unpruned. Per certificate it re-derives the Pell solution and its
minimality; the automorph, $P=X_0^{\,j}$, $P\equiv I\pmod M$, $Q$-invariance and $\gcd$
preservation; $B_{\mathrm{cut}}$ recomputed from Theorem~\ref{thm:cutoff} with the majorant
\eqref{eq:maj}; $\epsw<\epsbar$ and the blocker bound; the two cases of
Lemma~\ref{lem:backward}; $L^{*}$ with both witnesses \eqref{eq:witness} and $Y^{*}\ge Y_{\min}$;
$T_i>L_i^{*}$; the orbit list by an \emph{independent} scan of the reduction box; every ledger
row; $d^{*}$, $p_{\mathrm{noblock}}$ and $p_{\mathrm{cert}}$; the realised offsets; the
sufficiency statement of Lemma~\ref{lem:binding} against exhaustive scans; the tail verdicts
unpruned; each declared tie as an exact equality at lower height; the partition of
Proposition~\ref{prop:partition} on a bounded population; the completeness of $P_1$ and $P_5$ by
independent enumeration and set equality; and an independent brute-force scan.

\paragraph{Results obtained here.}
The archive is self-contained: running \texttt{emit.py} inside it regenerates all $26$
certificates \emph{byte-identically} (verified by \texttt{sha256} before and after), and running
\texttt{check.py} inside it reproduces the audit below exactly. The runs recorded in
\texttt{results/} are:
\begin{center}\small
\begin{tabular}{@{}lll@{}}
\toprule
run & result & log \\
\midrule
\texttt{check.py}, all 26 certificates & $1707/1707$ exact checks, $0$ failures &
  \texttt{check\_v8.log} \\
\quad the same, re-run inside the archive & $1707/1707$, $0$ failures &
  \texttt{check\_archive.log} \\
\texttt{emit.py} inside the archive & $26$ certificates, byte-identical &
  \texttt{emit\_archive.log} \\
\texttt{qr\_general.py} & $98/98$, $0$ failures & \texttt{qr\_general2.log} \\
\texttt{lit\_bridge.py} & $14/14$, $0$ failures & \texttt{lit\_bridge.log} \\
\texttt{review3.py} & $25/25$, $0$ failures & \texttt{review3.log} \\
\texttt{review4.py} & $14/14$, $0$ failures & \texttt{review4.log} \\
\texttt{review5.py} & $14/14$, $0$ failures & \texttt{review5.log} \\
\texttt{review6\_sail.py} & $14/14$, $0$ failures & \texttt{review6\_sail.log} \\
\bottomrule
\end{tabular}
\end{center}

\paragraph{What is not claimed.}
No part of this note is machine-checked in a proof assistant; none of it has been refereed; and
testing $26$ instances does not and cannot establish Theorem~\ref{thm:main}, which rests
on the proofs of Sections~\ref{sec:automorph}--\ref{sec:partition}.

\section{Limitations and extension questions}\label{sec:limits}

\paragraph{Limitations.}
\begin{enumerate}\itemsep2pt
\item \textbf{Cost.} Theorem~\ref{thm:main} asserts termination, not efficiency. No closed-form
bound on the running time is proved, and the dominant cost grows with $\rho=\lambda^{\,j}$ where
$j\le|GL_2(\Z/M\Z)|$ (Remark~\ref{rem:cost}). The constant
$d^{*}$ is computed per instance from $\rho$ and $B_{\mathrm{comp}}$; no single closed form
uniform in the input is written out.
\item \textbf{The finite region is scanned.} $P_1$ is decided by a direct scan of
\eqref{eq:P1box} rather than by the orbit machinery. That is legitimate for an effectiveness
theorem, and the region is small in every instance tested, but it is a scan.
\item \textbf{Weak records.} The theorem classifies \emph{strict} records. Ties are detected and
the certified tie families are genuine infinite families of weak records, but a full
classification of weak records --- in particular of imprimitive ones --- is not claimed and
would need a separate argument.
\item \textbf{$B_{\mathrm{cut}}$ is an upper bound}, not the supremum of \eqref{eq:cut}
(Warning~\ref{warn:sup}).
\item \textbf{Equality in Proposition~\ref{prop:classical} is open.} The inclusion is proved and
its strictness is witnessed; whether the two sets ever coincide as infinite sets is not decided.
The agreement observed at $\alpha=\sqrt3$ with $m,k$ odd holds over a finite range of heights
only and is recorded as an observation, not a theorem (Remark~\ref{rem:noteq}).
\item \textbf{Novelty.} Whether this objective, or a classification of its records, appears in
the literature is unsettled. Targeted searches found nothing matching, which is not evidence.
The objective is a weighted multiplicative-metric best-approximation problem on a union of
residue cosets, and such problems are a recognised family, so a prior treatment is plausible and
has not been excluded. The effectiveness, generality and certificate-completeness claims of
Section~\ref{sec:intro} are therefore labelled a \emph{proposed} contribution. The nearest
works consulted, and the specific distinctions, are these.
\begin{itemize}\itemsep2pt
\item \cite{KushwahaSarma} develops best approximation for
$\mathcal F_{p^\ell}$-continued fractions. The congruence structure there is supplied by a fixed
prime power and the approximants are reduced fractions compared under denominator ordering; here
$M$ and $\mathcal R$ are arbitrary, imprimitive pairs are admissible and are shown to occur
among records, and the order is $H$.
\item \cite{HanLee2025} proves moment formulas for Siegel transforms under congruence
conditions in dimension $2$, with a Schmidt-type counting theorem and a quantitative Khintchine
theorem for approximation by $p/q$ with congruence constraints on the pair. Those are counting
and metric statements --- almost-everywhere and on average, for primitive pairs --- whereas the
present question is the exact identification of individual record states for a fixed $\alpha$.
\item \cite{SpaldingVeselov} relates the Conway river of a real indefinite binary quadratic form
to the Arnold sail of the corresponding pair of lines; the sails there are attached to a
specified pair of lines. The forms are the same objects as $Q$ here, so the comparison is
natural, but Proposition~\ref{prop:sail} shows the records of $\varepsilon$ are not in general
hull vertices of their admissible cone, so no reading-off of records from such a sail is
available.
\end{itemize}
None of this establishes priority in either direction, and references that bear on the question
are welcome.
\item \textbf{$H$-order versus denominator order.} Proposition~\ref{prop:classical} compares two
notions both ordered by $H$. Converting to the denominator-ordered notion of
\cite{KimLeeLiao2022,KimLeeLiao2025,KushwahaSarma} is an unresolved bridge, and nothing here may
be read as a statement about best $1$-rational approximations in their sense. In particular the
witness $(7,11)$ of Proposition~\ref{prop:classical} is a statement about the $H$-ordered
comparison only and does not bear on the denominator-ordered one.
\item \textbf{$\mathcal R$ is a union of cosets}, not in general a sublattice of $\Z^2$; it is a
sublattice only when $\mathcal R$ is a subgroup containing $(0,0)$. The sail picture of
\cite{Karpenkov2013,SpaldingVeselov} therefore applies coset by coset at best, and
Proposition~\ref{prop:sail} rules out the natural coset-hull reduction outright. It does not rule
out some other geometric description of the record set, and finding one is an open question
rather than a closed door.
\item \textbf{Scope.} The method is specific to quadratic $\alpha$; it says nothing about
$\alpha$ of higher degree, and nothing in this note bears on any other conjecture.
\end{enumerate}

\paragraph{Extension questions.} Clearly separated from the theorem:
\begin{enumerate}\itemsep2pt
\item \emph{Sharper cutoffs and cost.} Is there a bound on $B_{\mathrm{cut}}$, or on the total
work, polynomial in $(\log\Delta,\ M)$? The obstruction identified here is
$B_{\mathrm{cut}}\ge\max_A|Q_A|$ over any family (Remark~\ref{rem:families}); is there a
family-free cutoff?
\item \emph{A bridge to denominator-ordered approximation.} Is there a transfer, in either
direction, between records of $\varepsilon$ ordered by $H$ and best approximations ordered by
the denominator, under a congruence restriction? Even a partial dictionary would connect this
classification to \cite{KimLeeLiao2022,KimLeeLiao2025,Elsner2000}.
\item \emph{Weak records.} Classify all weak records, including imprimitive ones. The certified
ties here suggest the equality mechanism is governed by the degenerate case $\Gamma_d=0$ of
Lemma~\ref{lem:epscmp}; is that the only source?
\item \emph{Other error objectives.} For which alternative objectives can one establish
analogues of the defect cutoff (Theorem~\ref{thm:cutoff}) and the eventual comparison lemmas
(Section~\ref{sec:compare})? Nothing here establishes that being a ratio of two linear forms in
$(m,k)$ is by itself sufficient; the cutoff proof uses the specific identities of
Lemma~\ref{lem:ident}.
\item \emph{Higher degree.} Is there any analogue when $\alpha$ is cubic? This is a
\emph{question}, not a transfer: the whole mechanism rests on a rank-one unit group acting on a
binary form, and no higher-degree substitute is proposed here.
\item \emph{A Lagrange-type spectrum.} Along a fixed-defect orbit
$H^2\varepsilon\to2\alpha|Q|/\sqrt\Delta$, so the accumulation values at records are explicit.
What is the envelope of $H\mapsto H^2\varepsilon$ \emph{between} records?
\end{enumerate}

\begin{thebibliography}{99}
\bibitem{BuellBQF} D.~A.~Buell. \emph{Binary Quadratic Forms: Classical Theory and Modern
Computations}. Springer, 1989. \href{https://doi.org/10.1007/978-1-4612-4542-1}
{doi:10.1007/978-1-4612-4542-1}.

\bibitem{ChaKim2023} B.~Cha and D.~H.~Kim. \emph{Intrinsic Diophantine approximation on the unit
circle and its Lagrange spectrum}. Ann. Inst. Fourier (Grenoble) \textbf{73} (2023), no.~1,
101--161. \href{https://doi.org/10.5802/aif.3522}{doi:10.5802/aif.3522}.

\bibitem{Cohen1993} H.~Cohen. \emph{A Course in Computational Algebraic Number Theory}. GTM
138, Springer, 1993. \href{https://doi.org/10.1007/978-3-662-02945-9}
{doi:10.1007/978-3-662-02945-9}.

\bibitem{ConradGenPell} \emph{Generalized Pell equation}. Math 154 course handout, Stanford
University. Theorem 1.1 (finiteness of the classes of solutions of $x^2-dy^2=n$ under the
norm-one units) and Theorem 2.1 (an explicit bounded region containing a representative of each
class). The handout carries no printed author line and is cited by title, course and URL.
\url{https://virtualmath1.stanford.edu/~conrad/154Page/handouts/genpell.pdf}.

\bibitem{DeJesus2026baseline} E.~De Jes\'us. \emph{Boundary and Interior Record Chronologies for
the Two-Mode Oscillator Spectrum}. Zenodo, 2026. \href{https://doi.org/10.5281/zenodo.21796640}
{doi:10.5281/zenodo.21796640} (version); concept DOI 10.5281/zenodo.21796540.

\bibitem{DeJesus2026oddodd} E.~De Jes\'us. \emph{Odd/Odd Records for a Relative Approximation
Error: Three Quadratic Classifications and an Exact Certificate}. Zenodo, 2026.
\href{https://doi.org/10.5281/zenodo.24046651}{doi:10.5281/zenodo.24046651} (version); concept
DOI 10.5281/zenodo.24046650.

\bibitem{Elsner2000} C.~Elsner. \emph{On Diophantine approximations with rationals restricted by
arithmetical conditions}. Fibonacci Quart. \textbf{38} (2000), no.~1, 25--34.

\bibitem{HanLee2025} J.~Han and S.~B.~Lee. \emph{Moment formulas of Siegel transforms with
congruence conditions in dimension 2}. Preprint, arXiv:2507.05905 [math.NT], 2025.
\href{https://arxiv.org/abs/2507.05905}{arXiv:2507.05905}.

\bibitem{Karpenkov2013} O.~Karpenkov. \emph{Geometry of Continued Fractions}. Algorithms and
Computation in Mathematics 26, Springer, 2013.
\href{https://doi.org/10.1007/978-3-642-39368-6}{doi:10.1007/978-3-642-39368-6}.

\bibitem{KimLeeLiao2022} D.~H.~Kim, S.~B.~Lee and L.~Liao. \emph{Odd-odd continued fraction
algorithm}. Monatsh. Math. \textbf{198} (2022), 323--344.
\href{https://doi.org/10.1007/s00605-022-01704-2}{doi:10.1007/s00605-022-01704-2}.

\bibitem{KimLeeLiao2025} D.~H.~Kim, S.~B.~Lee and L.~Liao. \emph{Diophantine approximation by
rational numbers of certain parity types}. Mathematika \textbf{71} (2025), e12285.
\href{https://doi.org/10.1112/mtk.12285}{doi:10.1112/mtk.12285}.
\bibitem{KushwahaSarma} S.~Kushwaha and R.~Sarma. \emph{Best approximations by
$\mathcal F_{p^\ell}$-continued fractions}. Preprint, arXiv:2112.00367 [math.NT], 2021.
\href{https://arxiv.org/abs/2112.00367}{arXiv:2112.00367}.

\bibitem{SpaldingVeselov} K.~Spalding and A.~P.~Veselov. \emph{Conway river and Arnold sail}.
Arnold Math. J. \textbf{4} (2018), 169--177. Preprint arXiv:1801.10072 [math.NT].
\href{https://arxiv.org/abs/1801.10072}{arXiv:1801.10072}.
\end{thebibliography}

\appendix

\section{Proof-to-certificate crosswalk}\label{app:crosswalk}

Each row names an obligation, where it is proved, the certificate fields that carry it, and the
check that re-derives it from the definitions. Check labels are those printed by
\texttt{check.py}.

{\footnotesize
\setlength{\tabcolsep}{3pt}
\begin{longtable}{@{}>{\raggedright\arraybackslash}p{5.0cm}
                    >{\raggedright\arraybackslash}p{2.3cm}
                    >{\raggedright\arraybackslash}p{5.2cm}
                    >{\raggedright\arraybackslash}p{1.1cm}@{}}
\toprule
obligation & proved in & certificate fields & check \\
\midrule
\endhead
$T^2-\Delta U^2=4$, $U$ minimal, $\Delta$ nonsquare, primitive, $\alpha>0$ &
  Lem.~\ref{lem:automorph} & \texttt{T}, \texttt{U}, \texttt{A,B,C,s}, \texttt{Delta} & C1 \\
$X_0$ proper automorph; $\det=1$; $Q$-invariant; $\eta\mapsto\lambda\eta$, $\lambda>1$ &
  Lem.~\ref{lem:automorph} & \texttt{X0} & C2 \\
$P=X_0^{\,j}\equiv I\!\!\pmod M$; preserves $Q$, classes, $\gcd$; $\eta\mapsto\rho\eta$ &
  Lem.~\ref{lem:automorph} & \texttt{P}, \texttt{j} & C3 \\
comparison family admissible, $|Q|\le D$, declared floor lowest &
  Lem.~\ref{lem:fam} & \texttt{family}, \texttt{D}, \texttt{family\_floor} & C4 \\
$B_{\mathrm{cut}}$ from \eqref{eq:cut} with the majorant \eqref{eq:maj} &
  Thm.~\ref{thm:cutoff}, Lem.~\ref{lem:tailbound} &
  \texttt{Bcut}, \texttt{maxF\_pair}, \texttt{cutoff\_periods}, \texttt{cutoff\_tail\_from} & C5 \\
$\epsw<\epsbar$; blocker bound non-vacuous; $B_{\mathrm{comp}}\le3B_{\mathrm{cut}}$ &
  Lem.~\ref{lem:G}, \ref{lem:blocker} &
  \texttt{eps0\_state}, \texttt{eps0\_num}, \texttt{Bcomp}, \texttt{Bcomp\_uniform} & C6 \\
$P_3$: no forward tail $\Rightarrow G<0$ and the floor holds &
  Lem.~\ref{lem:backward}(1) & per-orbit \texttt{case}$=$\texttt{i} & C6b \\
$P_4$: $Y\le Y^{*}\!\iff\!\varepsilon\ge\epsw$ on $G>0$ &
  Lem.~\ref{lem:backward}(2) & per-orbit \texttt{case}$=$\texttt{ii}, \texttt{Ystar\_num} & C6b \\
$L^{*}$ with both witnesses \eqref{eq:witness}; $Y_L$ strictly increasing &
  Lem.~\ref{lem:cutlevel} & per-orbit \texttt{Lstar}, \texttt{lvl\_eps0} & C6b \\
$Y^{*}\ge Y_{\min}$ when $BQ>0$ & Lem.~\ref{lem:ymin} & per-orbit \texttt{Ystar\_num},
  \texttt{Q} & C6b \\
$T_i=\mathrm{enter}_i+p_{\mathrm{cert}}$ and $T_i>L_i^{*}$ &
  \S\ref{sec:partition} & per-orbit \texttt{T}, \texttt{Lstar}, \texttt{enter} & C6b \\
enumeration complete to $B_{\mathrm{comp}}$; imprimitive retained; no PID input &
  Prop.~\ref{prop:box}, \ref{prop:split}, Lem.~\ref{lem:imprim} &
  \texttt{box}, \texttt{nclass}, \texttt{orbits}, \texttt{W} & C7 \\
every ledger row: $\Gamma_d$ sign, right-hand side, both thresholds &
  Lem.~\ref{lem:epscmp}, \ref{lem:hcmp} & \texttt{ledger} & C8 \\
$\varepsilon$ decreasing and $H$ increasing along every candidate orbit &
  Lem.~\ref{lem:epsmono}, \ref{lem:hcmp} & same-orbit $d{=}1$ rows & C8 \\
$d^{*}$, $p_{\mathrm{noblock}}$, and $p_{\mathrm{cert}}$ dominating both &
  Lem.~\ref{lem:dstar} & \texttt{d\_star}, \texttt{p\_noblock}, \texttt{p\_cert} & C9 \\
offset range complete at both ends: $d\le d^{*}$ and $d\ge0$ &
  Lem.~\ref{lem:dstar}, \ref{lem:lower} & \texttt{ledger} (both orders at $d{=}1$),
  \texttt{observed\_blocking\_offsets} & C9 \\
binding-state sufficiency against exhaustive scans & Lem.~\ref{lem:binding} &
  --- (behavioural) & C9 \\
$P_6$: tail verdicts unpruned; tails and ties reproduce; ties exact at lower height &
  Lem.~\ref{lem:epscmp}, \ref{lem:binding} & \texttt{record\_tails}, \texttt{ties} & C10 \\
exact ordering with a guaranteed exact-comparator fallback & Rem.~\ref{rem:sort} &
  --- (behavioural) & C5, C12 \\
the partition is exhaustive and each part's property holds &
  Prop.~\ref{prop:partition} & all of the above & C13 \\
$P_1$ complete from \eqref{eq:P1box}, record set recomputed &
  \S\ref{sec:partition} & \texttt{p1\_records}, \texttt{eps0\_state} & C14(a) \\
$P_5$ complete: every level $L_i^{*}<L<T_i$ enumerated, set equality &
  \S\ref{sec:partition}, Rem.~\ref{rem:completeness} & \texttt{initial\_exceptions},
  \texttt{Lstar}, \texttt{T} & C14(b) \\
\bottomrule
\end{longtable}}

One statement of the note is outside the certificate scheme, because it is a single finite
witness rather than a property of an instance: Proposition~\ref{prop:sail} is verified by
\texttt{scripts/review6\_sail.py} ($14/14$), which does not read the certificates.

\section{Certificate schema}\label{app:schema}

Each \texttt{certificates/NAME.json} is a flat JSON object. Its fields, grouped by the stage
that produces them:

\begin{center}\footnotesize
\setlength{\tabcolsep}{4pt}
\begin{tabular}{@{}>{\raggedright\arraybackslash}p{2.6cm}>{\raggedright\arraybackslash}p{11.3cm}@{}}
\toprule
stage & fields \\
\midrule
input &
\texttt{name}, \texttt{A}, \texttt{B}, \texttt{C}, \texttt{s}, \texttt{M}, \texttt{R},
\texttt{Delta}, \texttt{Dsqfree} \\
automorph &
\texttt{T}, \texttt{U}, \texttt{X0}, \texttt{j}, \texttt{P} \\
family, cutoff &
\texttt{family}, \texttt{D}, \texttt{family\_floor}, \texttt{Bcut}, \texttt{maxF\_pair},
\texttt{cutoff\_periods}, \texttt{cutoff\_tail\_from} \\
comparison state &
\texttt{eps0\_state}, \texttt{eps0\_num}, \texttt{Bcomp}, \texttt{Bcomp\_uniform},
\texttt{pruning}, \texttt{blocker\_fallback\_calls}, \texttt{verdict\_calls} \\
enumeration &
\texttt{box}, \texttt{nclass}, \texttt{W}, \texttt{window\_covers\_blockers} \\
per orbit &
\texttt{rep}, \texttt{Q}, \texttt{gcd}, \texttt{fwd}, \texttt{double}, \texttt{enter},
\texttt{win\_lo}, \texttt{win\_hi}, \texttt{case}, \texttt{Lstar}, \texttt{Ystar\_num},
\texttt{lvl\_eps0}, \texttt{T}, \texttt{tail} \\
levels, ledger &
\texttt{d\_star}, \texttt{p\_noblock}, \texttt{p\_cert}, \texttt{p\_cert\_binding},
\texttt{ledger} --- one entry per ordered pair and offset, carrying \texttt{A}, \texttt{B},
\texttt{d}, \texttt{QA}, \texttt{QB}, \texttt{rhs}, \texttt{Gd\_sign},
\texttt{eps\_threshold}, \texttt{eps\_verdict}, \texttt{h\_threshold},
\texttt{h\_verdict} \\
output &
\texttt{record\_tails}, \texttt{ties}, \texttt{initial\_exceptions}, \texttt{p1\_states},
\texttt{p1\_records} \\
diagnostics &
\texttt{observed\_binding\_offsets}, \texttt{observed\_binding\_offsets\_Gpos},
\texttt{observed\_blocking\_offsets} \\
\bottomrule
\end{tabular}
\end{center}

Exact algebraic quantities that are not integers --- the attained error $\epsw$ and the cutoff
$Y^{*}$ --- are stored as the four integers (numerator and denominator of each of the two
rational coordinates over $\sqrt{\texttt{Dsqfree}}$), so a checker reconstructs them without
rounding.

\end{document}
